% Exemplo: palavras cruzadas com o pacote cwpuzzle (bundle crossword), com termos do LaTeX. O mesmo
% arquivo gera a solução: veja palavras-cruzadas-solucao.tex. O QR code leva à solução no guia.
% Compilar: latexmk -pdf palavras-cruzadas.tex    Testado com: cwpuzzle 1.11, qrcode 1.51, TeX Live 2024 ou mais recente
%%BEGIN:cruzadaspreambulo
\documentclass[12pt, a4paper]{article}
\usepackage[T1]{fontenc}
\usepackage[brazilian]{babel}
\usepackage{cwpuzzle}                 % ambientes Puzzle e PuzzleClues
\usepackage{qrcode}
% a grade fica na mesma linha que o QR code (sem o \par antes e depois)
\renewcommand{\PuzzlePre}{}\renewcommand{\PuzzlePost}{}
\ifdefined\solucao \PuzzleSolution \fi   % a solução preenche as casas
%%END:cruzadaspreambulo

\begin{document}
\pagestyle{empty}

\begin{center}
%%BEGIN:cruzadasgrade
\PuzzleUnitlength=16pt               % lado de cada casa (o padrão é 20pt)
\begin{Puzzle}{14}{15}
  |{}   |[1]E|{}   |[2]I |{}|{}  |{}   |{}   |{}  |{}   |{}  |{}|[3]O|{}|.
  |{}   |X   |{}   |T    |{}|{}  |{}   |{}   |{}  |{}   |{}  |{}|V   |{}|.
  |[4]L |A   |B    |E    |L |{}  |{}   |{}   |{}  |{}   |{}  |{}|E   |{}|.
  |{}   |M   |{}   |M    |{}|{}  |{}   |{}   |[5]F|I    |G   |U |R   |E |.
  |{}   |{}  |[6]T |I    |K |Z   |{}   |{}   |{}  |{}   |{}  |{}|L   |{}|.
  |{}   |{}  |{}   |Z    |{}|{}  |{}   |{}   |{}  |{}   |[7]S|{}|E   |{}|.
  |[8]C |I   |T    |E    |{}|[9]B|{}   |[10]T|A   |B    |U   |L |A   |R |.
  |A    |{}  |{}   |{}   |{}|A   |{}   |{}   |{}  |{}   |B   |{}|F   |{}|.
  |[11]P|R   |E    |[12]A|M |B   |U    |[13]L|O   |{}   |F   |{}|{}  |{}|.
  |T    |{}  |{}   |C    |{}|E   |{}   |A    |{}  |[14]B|I   |B |E   |R |.
  |I    |{}  |{}   |R    |{}|L   |{}   |T    |{}  |{}   |L   |{}|{}  |{}|.
  |O    |{}  |[15]L|O    |G |{}  |[16]B|E    |A   |M    |E   |R |{}  |{}|.
  |N    |{}  |{}   |{}   |{}|{}  |{}   |X    |{}  |{}   |S   |{}|{}  |{}|.
  |{}   |{}  |{}   |{}   |{}|{}  |{}   |M    |{}  |{}   |{}  |{}|{}  |{}|.
  |{}   |{}  |{}   |{}   |{}|{}  |{}   |K    |{}  |{}   |{}  |{}|{}  |{}|.
\end{Puzzle}
%%END:cruzadasgrade
\qquad
%%BEGIN:cruzadasqr
\parbox[b]{2cm}{\centering   % ao lado da grade, alinhado pela base
  \qrcode[height=1.8cm]{https://github.com/Smarzaro/latex-academico/blob/main/exemplos/diversao/solucao-palavras-cruzadas.md}\\[2pt]
  \footnotesize Solução no guia}
%%END:cruzadasqr
\end{center}

%%BEGIN:cruzadaspistas
% cada PuzzleClues é uma coluna, alinhada pelo topo (o % as deixa lado a lado)
{\footnotesize\raggedright\noindent
\renewcommand{\PuzzleClueWidth}{.48\textwidth}
\begin{PuzzleClues}{\textbf{Horizontais}}
  \Clue{4}{LABEL}{Comando que dá nome a um elemento para a referência cruzada}
  \Clue{5}{FIGURE}{Ambiente flutuante para ilustrações}
  \Clue{6}{TIKZ}{Pacote para desenhos e diagramas}
  \Clue{8}{CITE}{Comando básico de citação}
  \Clue{10}{TABULAR}{Ambiente que monta a grade de uma tabela}
  \Clue{11}{PREAMBULO}{Parte do arquivo antes do início do documento (sem acento)}
  \Clue{14}{BIBER}{Programa que processa as referências do \texttt{biblatex}}
  \Clue{15}{LOG}{Arquivo em que ficam os erros e os avisos da compilação}
  \Clue{16}{BEAMER}{Classe para apresentações de slides}
\end{PuzzleClues}%
\begin{PuzzleClues}{\textbf{Verticais}}
  \Clue{1}{EXAM}{Classe para provas e listas de exercícios}
  \Clue{2}{ITEMIZE}{Ambiente de lista com marcadores}
  \Clue{3}{OVERLEAF}{Editor de \LaTeX\ na \textit{web}}
  \Clue{7}{SUBFILES}{Pacote que permite compilar cada aula separadamente}
  \Clue{8}{CAPTION}{Comando que cria a legenda numerada de figuras e tabelas}
  \Clue{9}{BABEL}{Pacote que adapta o documento ao português}
  \Clue{12}{ACRO}{Pacote de siglas}
  \Clue{13}{LATEXMK}{Programa que repete as compilações até tudo ficar certo}
\end{PuzzleClues}\par}
%%END:cruzadaspistas

\end{document}
